% !TeX root = ../教材题库_全部条目_转写工作稿.tex
% Source: XJ-B1-C04-TH-P110-Y1359
% 第4章 / 4.2 指数函数 / 教材 p.110、111 / PDF 第117、118页
% OCR未校对：公式与文字均可能有误；未完成数学TeX转写；禁止直接发给学生
% 原题及相邻原解截图（核对用）：../题目原图/XJ-B1-C04-TH-P110-Y1359-01.webp, ../题目原图/XJ-B1-C04-TH-P110-Y1359-02.webp
% 本文件只在转写完成后作为题干使用；OCR 条目尚不是可靠题干。
\begin{bankitem}{p.110 思考（含紧邻说明）}{第4章 / 4.2 指数函数 / 教材 p.110、111 / PDF 第117、118页}{XJ-B1-C04-TH-P110-Y1359}{OCR未校对：公式与文字均可能有误；未完成数学TeX转写；禁止直接发给学生}
\begin{ocrdraft}
[教材p.110 / PDF第117页]
\par
一）的图象之间有什么关系？你能利用
\par
？思考当a>l时，函数3，一与3'
\par
函数的图象画出函数
\par
的图象吗？
\par
1
\par
=a-x，因此当函数3一）和3，一a'的自变量取互为相反数的两个
\par
值时，对应的函数值相等．也就是说，如果点。，(o)在3一）的图象上，那么
\par
[教材p.111 / PDF第118页]
\par
这个点关于y轴的对称点（一“，\ocrglyph{2116}）一定在3，一的图象上；反之，3-的图象上
\par
任意一点（“，\ocrglyph{2116}），其关于y轴的对称点（亠。，！y。）也一定在3一）的图象上．
\par
因此，指数函数3一“。和y=(7)的图象关于y轴对称，如图4，2一4．
\par
y=d(a>l)
\par
底数互为倒数的两个
\par
指数函数的图象关于3，轴
\par
对称，
\par
0
\par
图4．2一4
\par
在图4，2一4中，若令上一乃，则0<b<1，于是y=（一
\par
>1）的图象即为3-
\par
俨（0<b<D的图象．
\par
综上可得，指数函数3-a'伍>0且磊乒1）的图象及性质如下：
\par
表达式
\par
y=aX(O<a<1)
\par
图象
\par
（0，l)
\par
（0，l)
\par
0
\par
0
\par
定义域
\par
值域
\par
函数图象过定点（0，1)，即=1
\par
在R上递减
\par
在R上递增
\end{ocrdraft}
\end{bankitem}
